% Generated by `taskgen course examples` from powmod.sliding_window (seed 6).
% Do not edit: change examples.yaml and regenerate.
\begin{example}[метод скользящего окна, $k = 3$]\label{ex:exp-sliding}
Вычислить \[293^{970}\ (\textrm{mod}\ 412)\] методом скользящего окна ширины $k = 3$ (предвычисляются нечётные степени $a, a^3, a^5, a^7$). Подсчитать число возведений в квадрат и умножений (вместе с предвычислением) и сравнить с бинарным алгоритмом.

\par\noindent\textit{Решение.} $e = 970 = 1111001010_2$. Предвычисление: $a^2 = 153$, $a^{3} = 333$, $a^{5} = 273$, $a^{7} = 157$ (возведений в квадрат: $1$, умножений: $3$). \par Биты просматриваются слева направо: нулевой бит --- окно <<$0$>> (одно возведение в квадрат); единичный бит открывает окно длины не больше $3$, оканчивающееся единицей; окно $w$ длины $s$: $r \gets r^{2^s} \cdot a^{w}$ (первое окно: $r \gets a^{w}$). Окна: \[111\,|\,1\,|\,0\,|\,0\,|\,101\,|\,0\]\[\begin{array}{c|l|c|r|r}\text{окно} & r^2, r^4, \dots & \text{множитель} & r & e \\ \hline 111 & \text{---} & a^{7} & 157 & 7 \\ 1 & 341 & a & 209 & 15 \\ 0 & 9 & \text{---} & 9 & 30 \\ 0 & 81 & \text{---} & 81 & 60 \\ 101 & 381,\ 137,\ 229 & a^{5} & 305 & 485 \\ 0 & 325 & \text{---} & 325 & 970 \\ \end{array}\]Число операций (вместе с предвычислением):\[\begin{array}{l|c|c|c} & \text{квадраты} & \text{умножения} & \text{всего} \\ \hline \text{окно } k = 3 & 8 & 5 & 13 \\ \text{бинарный} & 9 & 5 & 14 \\ \end{array}\]Экономия метода окна: $1$.\[293^{970} \equiv 325\ (\textrm{mod}\ 412).\]
\end{example}
